% GENERATED by tools/make_books.py --bib. Do not edit: regenerate.

\begin{longtable}{>{\raggedright\arraybackslash}p{0.30\textwidth}>{\raggedright\arraybackslash}p{0.46\textwidth}>{\raggedright\arraybackslash}p{0.16\textwidth}}
\toprule \textbf{Source} & \textbf{What it implements} & \textbf{References} \\ \midrule \endfirsthead
\toprule \textbf{Source} & \textbf{What it implements} & \textbf{References} \\ \midrule \endhead
\bottomrule \endfoot
\code{src/research/research.cpp} & Taubin λ\textbackslash{}|μ smoothing (\code{taubin\_smoothing}) (implemented (example)) & \cite{Taubin1995Signal} \\
\code{src/research/research.cpp} & Add Surface Noise (\code{surface\_noise}), noisy input for comparing smoothers (utility) & (utility, no paper) \\
\code{docs/tasks/0235-hci-pass.md}, \code{src/editor/tab\_strip.cpp} & The HCI audit of the editor's windows, tabs, menus and bars (task 0235): visibility of status (tabs never clipped, the breadcrumb), recognition over recall (tab groups, a stage's tabs), consistency (one signal per tab state), memory load (the guide hides after two sessions), help (a tool's next step in the status bar) (applied (an audit, not an algorithm)) & \cite{Nielsen1994Enhancing} \\
\code{docs/tasks/0235-hci-pass.md} & Reach on a 13-inch 2-in-1: related tabs folded into one with a member row (one target instead of three), the hidden tabs one button away (+N), the Inspector folded with one button (task 0235) (applied) & \cite{Fitts1954Information} \\
\code{src/scene/mesh.cpp} & Catmull-Clark subdivision (core Mesh tool + SubdivisionSurface component) (implemented (core)) & \cite{Catmull1978Recursively} \\
\code{src/render/voxel\_gi.cpp} & Realtime GI (Lighting window > Realtime GI): binary voxel columns, the mip-map ray test (sec. 4), one bounce of the sun through a reflective shadow map and sky occlusion (sec. 5.1) in the rasterized views, with baked lightmaps (task 0013, ADR 0011) (implemented (renderer, CPU)) & \cite{Thiedemann2011Voxel} \\
\code{src/editor/sketch\_build.cpp} & Solids from a few strokes: a box from three lines at a corner, a cylinder or cone from two circles, a sweep from two circles and a path (\code{sketchbuild}); here on the sketch's scene canvases instead of a photograph, the solids as F9 operators (ADR 0065) (implemented (task 0225)) & \cite{Chen2013Sweep} \\
\code{src/research/sketch\_lift.cpp} & A scaffold of construction lines lifted from the picture plane: shared junctions, lines along the axes' vanishing directions, depths by least squares (\code{sketchlift}, ADR 0067); the symmetry and curved lines of the papers not taken (experimental, gate met (task 0227)) & \cite{Gryaditskaya2020Lifting} \\
\code{src/scene/sketch.cpp} & Pen strokes corrected at the lift: split at corners, fitted to lines, arcs, ellipses and curves, pulled to the guides (ADR 0063) (implemented (task 0223)) & \cite{Igarashi1997Interactive} \\
\code{src/scene/dynamesh.cpp} & DynaMesh's inside test: the winding number along grid lines of three axes, the nonzero rule merging self-intersections, a majority of three votes a grid vertex (ADR 0076) (implemented (the grid-line form, not the paper's hierarchical evaluation)) & \cite{Jacobson2013Robust} \\
\code{src/scene/dynamesh.cpp} & DynaMesh's surface: naive surface nets over the sign grid (a vertex per mixed cell at the mean of its edge crossings, a quad per mixed grid edge), then tangential relaxation (implemented) & \cite{Gibson1998Constrained} \\
\code{src/scene/quad\_remesh.cpp} & Quad Remesh: Blender's patched QuadriFlow compiled in (\code{extern/quadriflow}), driven as Blender's wrapper does, on a worker with a cancel flag (ADR 0077) (integrated (library)) & \cite{Huang2018Quadriflow} \\
\code{src/scene/dyntopo.cpp} & Dyntopo: edges under the brush longer than the detail size split at their midpoints, shorter ones collapsed where that folds no triangle (ADR 0078) (implemented) & (no paper: Sculptris / Blender's Dyntopo) \\
\code{src/render/mesh\_maps.cpp} & Mesh maps baked into the UV layout: world normal, position in the bounds, a sagitta-over-chord curvature smoothed twice, cosine-weighted occlusion and thickness through the path tracer on the mesh alone (ADR 0079) (implemented) & (no paper: Substance Painter's mesh maps) \\
\code{src/texture/surface.cpp}, \code{src/texture/smart\_material.cpp} & The surface's stack with mask stacks reading the mesh maps, the rect engine redoing only the edited layers inside the edited rectangle, smart materials as saved groups, texture sets packed per engine (ADRs 0080, 0082) (implemented) & (no paper: Substance Painter's layer stack, smart masks and smart materials) \\
\code{src/editor/paint.cpp} & Painting by projection: the brush's disc projected through the camera, texels placed by barycentrics and depth-tested against the view's planes; a stroke one undo step of the 64 x 64 tiles it touched (ADR 0081) (implemented) & (no paper: Blender's Project Paint) \\
\code{src/editor/paint.cpp} & The painterly brush's wet mixing: the brush carries the mean colour it went over and mixes the next dab by a wet amount - a one-value pickup, not IMPaSTo's paint volume and Kubelka-Munk mixing (ADR 0083) (simplified) & \cite{Baxter2004Impasto} \\
\code{src/scene/stroke\_object.cpp} & Stroke objects' Canopy normals: each card's custom normal points away from a centre so a bush of flat cards shades as one round volume (ADR 0084) (implemented) & (no paper: the hand-painted foliage trick) \\
\code{docs/hci/klm.md}, \code{tools/hci\_klm.py} & Nineteen frequent actions counted before and after the HCI pass (K, P, B, H, M with the paper's rules for M), recounted by a tool (ADR 0085) (applied (a model, not an algorithm)) & \cite{Card1980Keystroke} \\
\code{docs/hci/walkthrough-*.md} & Four pipelines walked at the baseline and after the pass, the four questions per step (ADR 0085) (applied) & \cite{Polson1992Cognitive} \\
\code{src/editor/sculpt.cpp} & Sculpt polygroups in the eight colours, then lighter and darker tiers (task 0269) (implemented) & \cite{Okabe2008Color} \\
\code{src/scene/export.cpp} & Skinned meshes exported with joints, inverse bind matrices, JOINTS\_0 / WEIGHTS\_0 and sampled channels (task 0268) (implemented) & \cite{Khronos2021Gltf} \\
\code{tools/posemodel/train.py}, \code{src/scene/puppet.cpp} & Puppet Posing's pose prior: any subset of the ends' goals to the whole pose, a residual network trained with goals hidden (task 0280) (implemented (much smaller, positions only)) & \cite{Oreshkin2022Protores} \\
\code{src/editor/ui.cpp} & The relative-luminance contrast ratio (\code{ui::contrast\_ratio}) measuring every theme; the High Contrast theme at 4.5:1 or more (task 0261) (implemented) & \cite{W3C2018Web} \\
\code{docs/hci/study-protocol.md}, \code{docs/hci/study-282.md} & The study's questionnaire (ADR 0085) (protocol written, **not run**) & \cite{Brooke1996Sus} \\
ufbx v0.20.0 & FBX import & \cite{libufbx} \\
fast\_float 5.0.0 & fast, exact float parsing in the scene loader & \cite{libfastfloat} \\
MikkTSpace & tangents for normal maps & \cite{libMikkTSpace} \\
Hosek-Wilkie sky model & physical sky & \cite{libHosekWilkieskymodel} \\
QuadriFlow 27a6867 & quad remeshing (task 0241); with it LEMON 1.3.1 (Boost Software License 1.0) and pcg32 (Apache-2.0); pss (Intel, BSD-3-Clause) is in the copy but not compiled & \cite{libQuadriFlow} \\
cuew & loads NVIDIA's CUDA driver at run time (builds with OptiX: the Windows and Linux installers since task 0169) & \cite{libcuew} \\
NVIDIA OptiX's loader & compiled into builds with OptiX's headers - the Windows and Linux installers since task 0169 (ADR 0054; ADR 0037 before it) & \cite{libNVIDIAOptiXsloader} \\
The OptiX kernel's PTX & Blendity's own GPL kernel compiled with NVRTC against NVIDIA's OptiX device headers (their \code{\_optix\_*} intrinsic calls), committed (task 0169, ADR 0054) & \cite{libTheOptiXkernelsPTX} \\
Embree 4.4.1 & path tracer ray casting & \cite{libEmbree} \\
OpenImageDenoise 2.5.0 & path tracer denoising & \cite{libOpenImageDenoise} \\
oneTBB 2022.3.0 & job system & \cite{liboneTBB} \\
Eigen (commit 8a1083e) & LSCM sparse solver (header-only) & \cite{libEigen} \\
OpenSubdiv 3.7.0 & optional Catmull-Clark backend & \cite{libOpenSubdiv} \\
OpenEXR + Imath 3.4.10 / 3.2.2 & EXR read/write & \cite{libOpenEXRImath} \\
OpenJPH 0.25.2 & runtime dependency of OpenEXR & \cite{libOpenJPH} \\
OpenColorIO 2.5.0 & AgX / Filmic / other view transforms & \cite{libOpenColorIO} \\
Manifold 3.5.2 & Boolean modifier & \cite{libManifold} \\
meshoptimizer 1.1 & Decimate modifier & \cite{libmeshoptimizer} \\
Open PGL 0.7.1 & path guiding & \cite{libOpenPGL} \\
Jolt Physics 5.6.0 & Play-mode rigid bodies & \cite{libJoltPhysics} \\
libjpeg-turbo 2.1.3 & JPEG decoding & \cite{liblibjpegturbo} \\
libpng 1.6.58 & PNG decoding & \cite{liblibpng} \\
zlib 1.3.1 & PNG compression & \cite{libzlib} \\
Zstandard 1.5.7 & compressed \code{.scene} files & \cite{libZstandard} \\
DPC++ runtime 7.1.0 & runtime dependency of Embree's build (SYCL) & \cite{libDPCruntime} \\
shaderc v2025.4 & compiles the GPU path-tracing kernel (GLSL to SPIR-V) at run time & \cite{libshaderc} \\
Vulkan-Headers 1.4.341 & GPU rendering API declarations (headers only; the Vulkan loader comes with the GPU driver) & \cite{libVulkanHeaders} \\
OpenXR SDK 1.1.65 & VR View (round 383, ADR 0096): the loader, loaded at run time, reaches the headset's OpenXR runtime & \cite{libOpenXRSDK} \\
FFmpeg 8.1 & Render Animation's MP4, encoded in-process (task 0077) & \cite{libFFmpeg} \\
x264 35fe20d & H.264 encoding, inside FFmpeg & \cite{libx264} \\
x265 4.1 & H.265 encoding, inside FFmpeg (Linux, macOS) & \cite{libx265} \\
libaom 3.13.1 & AV1, inside FFmpeg (\code{aom.dll} on Windows) & \cite{liblibaom} \\
libvpx, Opus, libvorbis / libogg / libtheora 1.15.2, 1.3.1, - & codecs inside FFmpeg & \cite{liblibvpxOpuslibvorbislibogglibth} \\
LAME 3.100 & MP3, inside FFmpeg & \cite{libLAME} \\
\end{longtable}
